\documentclass[11pt,a4paper]{article}
\usepackage[T1]{fontenc}
\usepackage{lmodern,amsmath,amssymb,amsthm,mathtools,booktabs,microtype}
\usepackage[margin=27mm]{geometry}
\usepackage[colorlinks=true,linkcolor=blue,citecolor=blue,urlcolor=blue]{hyperref}
\setlength{\emergencystretch}{2em}
\newtheorem{theorem}{Theorem}[section]
\newtheorem{lemma}[theorem]{Lemma}
\newtheorem{proposition}[theorem]{Proposition}
\newtheorem{definition}[theorem]{Definition}
\newcommand{\Q}{\mathbb Q}\newcommand{\Z}{\mathbb Z}
\DeclareMathOperator{\Tr}{Tr}
\DeclareMathOperator{\Nm}{N}
\DeclareMathOperator{\rad}{rad}
\title{A single-cover criterion from a mixed-power height inequality}
\author{}\date{17 November 2026}
\begin{document}\maketitle

\begin{abstract}
We prove a height inequality for a rational sum of an even power and a
square under a weak rational \(abc\) inequality. Applied to the elliptic
index construction, it places the required height information in one
detector argument while retaining one free rational parameter. A single
native rational-square condition then gives soundness for sufficiently
large fixed coefficient shifts. Its finite algebra has rank \(128\),
and the conditional polynomial census is nine variables, including the
input. A defining polynomial with that census additionally requires a
rational realization theorem for the actual single cover. That realization theorem remains open.
\end{abstract}

\section{A mixed-power height estimate}

For \(a/b\in\Q\) in lowest terms with \(b>0\), set
\(h(a/b)=\log\max(|a|,b)\).
Fix \(\kappa>1\) and \(K_\kappa\ge1\), and assume
\begin{equation}\label{eq:abc}
 c\le K_\kappa\rad(abc)^\kappa
\end{equation}
for every positive coprime integer triple \(a+b=c\).

\begin{lemma}\label{lem:mixed}
Let \(d\) be a positive even integer such that
\[
 \epsilon=d-\kappa(d/2+2)>0.
\]
For \(a,z\in\Q\) with \(u=a^d+z^2>0\), one has
\[
 h(a)\le C_{\kappa,d}h(u)+\epsilon^{-1}\log K_\kappa,
\]
where
\begin{equation}\label{eq:C}
 C_{\kappa,d}=
 \frac{\kappa(3d/2+1)}{\epsilon}+2
 =
 \frac{(4+\kappa)d-6\kappa}{(2-\kappa)d-4\kappa}.
\end{equation}
\end{lemma}
\begin{proof}
If \(a=0\), the assertion is immediate. If \(z=0\), then
\(h(a)=h(u)/d\), which is stronger. Assume both are nonzero.

Write \(u=C/F\) in lowest terms with \(C,F>0\). Multiplication by \(F^d\)
gives
\[
 (Fa)^d+(F^{d/2}z)^2=c',\qquad c'=CF^{d-1}\in\Z_{>0}.
\]
Write \(Fa=A/D\) in lowest terms, with \(D>0\).
The numerator of \(c'-(A/D)^d\) is coprime to \(D\).
Since \(d\) is even and this rational number is a square, there is a
nonzero integer \(B_0\) such that
\[
 F^{d/2}z=B_0/D^{d/2},\qquad A^d+B_0^2=c'D^d.
\]
In particular \(\gcd(A,D)=\gcd(B_0,D)=1\).
For \(g=\gcd(A^d,B_0^2)\), it follows that \(g\mid c'\).
After division by \(g\), the three terms are positive and pairwise
coprime. The weak inequality gives
\[
 \frac{c'D^d}{g}
 \le K_\kappa\rad(A B_0 c'D)^\kappa
 \le K_\kappa(|A|\,|B_0|\,c'D)^\kappa.
\]
Moreover,
\[
 |A|\le(c')^{1/d}D,\qquad
 |B_0|\le(c')^{1/2}D^{d/2}.
\]
Using \(c'/g\ge1\), we obtain
\[
 D^\epsilon\le K_\kappa(c')^{\kappa(3/2+1/d)}.
\]
Thus
\[
 h(Fa)\le\epsilon^{-1}\log K_\kappa+
 \left(\frac{\kappa(3/2+1/d)}{\epsilon}+\frac1d\right)\log c'.
\]
Finally, \(\log c'\le d\,h(u)\), \(\log F\le h(u)\), and
\(h(a)\le h(Fa)+\log F\). These inequalities give \eqref{eq:C}.
\end{proof}

\section{The single detector}

Retain the curve, point charts, subgroup identification and trace-height
estimate of \cite{Conditional}:
\[
 K=\Q(\sqrt2),\qquad E:y^2=x^3-(3-\sqrt2)^2x,\qquad
 2E(K)=\Z P,
\]
\[
 P=\left(\frac{57-28\sqrt2}{8},
          \frac{-660+417\sqrt2}{32}\right).
\]
In particular, a rational point-chart solution gives nonzero integers
\(n,h,j\) with
\[
 A=[n]P,\qquad H=[h]P,\qquad U=[j]P.
\]
There are fixed constants \(A_0>0,J_0\ge0\) such that
\begin{equation}\label{eq:traceheight}
 n^2\le A_0 h(T(n))+J_0,\qquad
 T(n)=\Tr_{K/\Q}x([n]P).
\end{equation}
The arithmetic hypotheses used for these identifications have the
scope stated in \cite{Conditional}.

Put \(f=1/x\), \(\zeta=-x/y\), and let a bar denote conjugation in \(K/\Q\).
Use the six symmetric tests of \cite{Reweighted}:
\begin{equation}\label{eq:tests}
 \begin{gathered}
 r=\frac{\zeta(U)}{\zeta(H)},\qquad
 s=\frac{\zeta(U-A)}{\zeta(H-P)},\\
 C_1=f(H),\quad C_2=f(U),\quad
 C_3=\overline{f(H-P)},\quad C_4=\overline{f(U-A)},\\
 C_5=r-\bar s,\qquad C_6=r+\bar s-2X.
 \end{gathered}
\end{equation}
Retain every existing point, difference and quotient guard.

Fix an odd integer \(B\ge3\), an even positive integer \(e\), and positive
integers \(d,m\), with \(d\) even. Set
\begin{equation}\label{eq:argument}
 \beta=X^2+B,\qquad T=\Tr x(A),\qquad
 H_0=\beta^e+T^e,\qquad t=H_0^{dm}+z^2,
\end{equation}
where \(z\in\Q\) is the sole free detector parameter.
There is no additional coordinate for \(t\) or \(H_0\).

Use the original fast detector of \cite{Reweighted}:
\begin{equation}\label{eq:detector}
 \begin{gathered}
 F(t)=\frac{t(t+1)(2t+1)(3t+1)(4t+1)(5t+2)}{120},
 \qquad w=\frac{F(t)}{4(t-1)^8},\\
 W=\prod_{i=1}^6\left(1+\frac{C_i^2}{F(t)}\right),\qquad
 V=\Tr(W^{-1}),\qquad J=\Nm(W^{-1}).
 \end{gathered}
\end{equation}
Guard \(F(t)\), \(t-1\), \(\Nm(W)\), \(V\), and all other displayed
denominators. The single finite square condition is
\begin{equation}\label{eq:square}
 {\cal R}=
 \left(1-\frac{V}{1+w}+\frac{J}{(1+w)^2}\right)
 \left(1+\frac{1-V/2}{w}\right)\in\Q^{\times2}.
\end{equation}
These are the actual detector trace and norm.

Let \(S\) be the fixed exceptional prime set of the conditional source,
and let \(\mathcal M(t)\) be the product of its nonexceptional odd
contact primes. The contact set is
\[
 \Theta=\{0,-1,-1/2,-1/3,-1/4,-2/5\}.
\]
For the parameters used below, \cite{Conditional} proves
\begin{equation}\label{eq:heights}
 \begin{gathered}
 h(t)\le\gamma^{-1}\log\mathcal M(t)+O(1),
       \qquad \gamma=8/\kappa-6>0,\\
 h_2(a,b)\le\alpha_e^{-1}h(a^e+b^e)+O(1),
       \qquad \alpha_e=e/\kappa-3>0,
 \end{gathered}
\end{equation}
where \(h_2(a,b)=h(1:a:b)\). The constants depend only on the fixed
parameters and \(K_\kappa\), not on \(B\) or a source tuple.

\section{Soundness}

\begin{lemma}\label{lem:local}
A guarded rational solution of \eqref{eq:square} has no negative odd
valuation of \(b=4w/V^2\) outside \(S\).
If \(X\ne n\), then
\[
 \mathcal M(t)\mid N,\qquad
 N=|\operatorname{num}(X-n)|.
\]
\end{lemma}
\begin{proof}
This detector is exactly the conditional detector of
\cite{Reweighted}; only the chosen argument has changed.
Every odd nonexceptional valuation of its physical weight \(w\)
is positive. At such a valuation, the reciprocal factor argument
puts each local slot of \(W^{-1}\) in the valuation ring, with residue
zero or one. If \(b\) had negative odd valuation, then \(V\) and \(J\)
would reduce to zero. The first factor of \({\cal R}\) would be a unit
with residue one, and the second would have negative odd valuation.
This contradicts \({\cal R}\in\Q^{\times2}\), as in \cite{Native}.

At an odd contact \(p\mid\mathcal M(t)\), write
\(\nu=v_p(F(t))=v_p(w)>0\), an odd integer. Each detector factor at
a place above \(p\) has valuation
\[
 \min(0,2v(C_i)-\nu)\le0.
\]
Since \(v_p(b)=\nu-2v_p(V)\) is not negative, there is a place above
\(p\) at which \(v(W^{-1})\le\nu/2\). Every summand in the valuation
of \(W\) is at least its total, so
\[
 v(C_i)>\nu/4>0\qquad(1\le i\le6)
\]
at that place. The strict inequality uses that \(\nu\) is odd.

The first four tests put \(H,U\) in the formal group at that place
and \(H-P,U-A\) in the formal group at the conjugate place.
The unchanged formal-logarithm estimates of \cite{Conditional}
compare \(r,\bar s\) with \(j/h,(j-n)/(h-1)\), with relative
error divisible by \(p^4\). At least one of \(h,h-1\) is a
\(p\)-adic unit, so the corresponding exact ratio is integral.
The congruence \(r-\bar s\equiv0\pmod p\) makes the other ratio
integral too, and the relative errors become absolute errors.
The symmetric last test then gives
\[
 r-X\equiv\bar s-X\equiv0\pmod p.
\]
Multiplying the two ratio congruences by \(h,h-1\), respectively,
and subtracting yields \(X-n\equiv0\pmod p\).
In particular \(X\) is \(p\)-integral. This proves the divisibility
at every contact prime.
\end{proof}

\begin{theorem}\label{thm:sound}
Suppose \eqref{eq:abc}, \eqref{eq:traceheight}, the retained point-chart
and local inputs hold. Assume
\[
 \epsilon>0,\quad \alpha_e>0,\quad \gamma>0,\qquad
 q_*=\frac{C_{\kappa,d}}{m\alpha_e\gamma}<2.
\]
There is a constant \(B_*\), independent of the source tuple, such that
every guarded rational solution of \eqref{eq:argument}--\eqref{eq:square}
with fixed odd \(B\ge B_*\) satisfies \(X=n\in\Z\).
\end{theorem}
\begin{proof}
Apply Lemma~\ref{lem:mixed} with \(a=H_0^m\).
Since \(h(H_0^m)=m h(H_0)\), it gives
\[
 h(H_0)\le\frac{C_{\kappa,d}}m h(t)+O(1).
\]
Combining \eqref{eq:heights} yields, for
\(E_h=h_2(\beta,T)\),
\[
 E_h\le q_*\log\mathcal M(t)+O(1).
\]
Every constant is independent of \(B\).

Suppose \(X\ne n\). Lemma~\ref{lem:local} gives
\(\mathcal M(t)\mid N\). If \(X=a/b\) is reduced, then
\[
 N=|a-nb|\le e^{h(X)}(1+|n|).
\]
The elementary identity \(\beta=X^2+B\), with \(B\ge3\), gives
\(2h(X)\le h(\beta)\le E_h\). Hence
\[
 E_h\le q_*\bigl(E_h/2+\log(1+|n|)\bigr)+O(1).
\]
Because \(q_*<2\), this bounds \(E_h\) by a fixed multiple of
\(\log(1+|n|)\) plus a fixed constant. Equation~\eqref{eq:traceheight}
then bounds \(n^2\) by the same expression. Thus \(n\) and \(E_h\)
are bounded independently of \(B\). But
\(E_h\ge h(\beta)\ge\log B\), a contradiction for every sufficiently
large \(B\).
\end{proof}

\begin{proposition}[Absorbing the outer amplification]\label{prop:amplification}
For fixed \(\kappa,d\), replacing \((e,m)\) by \((em,1)\)
preserves the total degree of the argument in \(\beta,T,z\)
and does not increase the soundness coefficient \(q_*\).
The improvement is strict if \(m>1\). For positive integral coefficients,
the replacement also decreases the height argument coefficientwise.
\end{proposition}
\begin{proof}
One has
\[
 \alpha_{em}=em/\kappa-3=m\alpha_e+3(m-1).
\]
The other two factors in \(q_*\) are unchanged, proving the
assertion about the soundness coefficient. Both height arguments have
highest degree \(edm\) in \(\beta,T\).
The terms of \((\beta^{em}+T^{em})^d\) are a subset of those of
\((\beta^e+T^e)^{dm}\). At the corresponding terms,
\(\binom d j\le\binom{dm}{mj}\): choosing \(j\) complete blocks of size
\(m\) embeds the former choices into the latter. Thus every coefficient
of the difference is nonnegative. This remains true after the
substitutions \(\beta=X^2+B\) and
\(T=2h_{A,0}^2+4h_{A,1}^2\), with \(B>0\).
\end{proof}

The following exact choices satisfy all strict inequalities:
\begin{center}
\begin{tabular}{rrrrrr}\toprule
\(\kappa\)&\(e\)&\(d\)&\(m\)&\(dm\)&\(q_*\)\\\midrule
\(25/24\)&8&6&1&6&\(10000/5187\)\\
\(16/15\)&10&6&1&6&\(32/17\)\\
\(121/100\)&16&12&1&12&\(40453083/21236816\)\\\bottomrule
\end{tabular}
\end{center}
Ordinary rational \(abc\) supplies the first row. The other rows retain
the indicated weaker inequalities. Their detector arguments are
\[
 (\beta^8+T^8)^6+z^2,\qquad
 (\beta^{10}+T^{10})^6+z^2,\qquad
 (\beta^{16}+T^{16})^{12}+z^2,
\]
respectively. In the retained input and point coordinates their total
degrees are \(96\), \(120\), and \(384\). These are degrees of the
height arguments, not determinant degree bounds. The argument does
not compute a numerical sufficient value of \(B_*\).

\section{Guards, finite encoding, and the remaining arithmetic assertion}

\begin{proposition}\label{prop:guards}
At a prepared actual point tuple on the inherited whole-character
guard locus, every new scalar denominator and the radicand
\({\cal R}\) are nonzero at every point-algebra character.
There is no nonzero condition on the retained parameter \(z\).
\end{proposition}
\begin{proof}
All such characters are real as in \cite{Reweighted,Native}.
Since \(e\) is even and \(\beta\ge B\ge3\),
\[
 H_0=\beta^e+T^e>0,\qquad
 t=H_0^{dm}+z^2>1.
\]
Consequently \(F(t)>0\), \(w>0\), and \(g=1+w>1\).
Both real slots of \(W\) are at least one. Thus their reciprocals
are positive and at most one, and \(0<V\le2\).
The first factor in \eqref{eq:square} is
\[
 \Nm(1-(gW)^{-1})>0,
\]
while the second is at least one. Hence \({\cal R}>0\).
The asserted denominator and norm guards follow.
The parameter \(z\) is unrestricted and appears through its square
only in \eqref{eq:argument}; \(z=0\) is allowed.
\end{proof}

\begin{definition}\label{def:realization}
For a fixed parameter row and fixed odd \(B\ge3\), the single-cover
realization property means that every nonzero integer \(X\) admits
actual source points with \(A=[X]P\) and a rational \(z\), on the
inherited whole-character point and difference guard locus, such
that the actual radicand \eqref{eq:square} is a nonzero rational square.
\end{definition}

\begin{theorem}\label{thm:count}
There is a computable nine-variable integer polynomial family whose
nonzero input projection is the guarded source above. Under the
hypotheses of Theorem~\ref{thm:sound}, and for every odd \(B\ge B_*\)
satisfying Definition~\ref{def:realization}, its input projection,
after adjoining zero by the usual input factor, is exactly \(\Z\).

The single-cover realization property is not proved in this paper.
Without it, Theorem~\ref{thm:sound} supplies the soundness direction
only; a nine-variable defining polynomial has not been established.
\end{theorem}
\begin{proof}
The three retained point pairs have the original finite point algebra
of rank \(4^3=64\). In that localized algebra, write
\({\cal R}=N_0/D_0\). Adjoin one monic quadratic generator
\[
 Y^2=N_0D_0,
\]
with the denominator and radicand norms in the guard product.
Proposition~\ref{prop:guards} proves that these added guards are units
at every prepared character. The total finite-algebra rank is
\(64\cdot2=128\).

The one-coordinate guarded finite-algebra encoder of \cite{Unit}
applies with the same descent and unit-guard hypotheses.
The unrestricted coordinates are the input, the six point coordinates,
\(z\), and the encoder coordinate: nine in total.
All rational expressions have finite computable integer coefficient
specifications after clearing denominators. No separate coordinate
or root is used for \(H_0\) or \(t\).

Soundness is Theorem~\ref{thm:sound}.
Definition~\ref{def:realization} gives a rational point character and
a rational root for the one added quadratic equation, proving
completeness. Multiplication by the input includes \(X=0\).
\end{proof}

\section{Fixed-place preparation and scope}

For fixed actual point data, as \(z\) tends to infinity at a fixed
place,
\[
 t\sim z^2,\qquad
 w\sim\frac1{4t^2},\qquad
 W-1=O(t^{-6})=O(w^3).
\]
The native radicand therefore satisfies
\[
 \frac{{\cal R}}{(w/(1+w))^2}\longrightarrow1.
\]
Nonzero local squares are open in every \(\Q_p^\times\).
A rational \(z\) can be chosen large simultaneously at finitely many
finite places and at the real place. This gives the same fixed-place
square preparation as \cite{Native}; the depths at \(2,3,17\) can
be made arbitrarily large. These statements concern a prescribed
finite set of places, not all primes.

For a fixed integer input and root point \(A=[X]P\), set
\(c=H_0^{dm}\). Writing \(z=a/b\), the contact numerators now come
from the six binary quadratic expressions
\[
 a^2+(c-\theta)b^2,\qquad \theta\in\Theta,
\]
with their fixed denominators cleared.
The inherited prime-pattern assertion with an unrestricted numerator
does not automatically give these representations or their required
conjugate reduction-order conditions. Even such a preparation theorem
would not prove that the actual radicand \({\cal R}\) is a rational square.

The remaining arithmetic problem is precisely a rational point on
this single actual double cover, allowing the auxiliary point indices
to vary and retaining every guard. The fixed-place limits and the
one-coordinate count do not solve it. No unrestricted twenty-one-variable
transfer follows from the mixed-power height estimate, which uses
the stated weak \(abc\) inequality.

\begin{thebibliography}{9}\raggedright
\bibitem{Conditional}
\emph{Thirteen Variables from Direct Elliptic Index Comparison},
repository manuscript,
\nolinkurl{papers/conditional_few_variables/direct_input13_20261110_v22.tex}.
\bibitem{Reweighted}
\emph{Quadratic reweightings of native elliptic detectors},
repository manuscript,
\nolinkurl{papers/pole_parameters/reweighted_elliptic_detectors_20261116_v53.tex}.
\bibitem{Native}
\emph{A guarded rational-square criterion for normalized elliptic detectors},
repository manuscript,
\nolinkurl{papers/pole_parameters/native_square_criterion_20261116_v53.tex}.
\bibitem{Unit}
\emph{A single norm determinant for the 25- and 13-variable constructions},
repository manuscript,
\nolinkurl{papers/finite_algebra/unit_norm_20261114_v28.tex}.
\end{thebibliography}
\end{document}
